\documentclass[10pt,a4paper]{amsart}
\usepackage[margin=19mm]{geometry}
\usepackage{amsmath,amssymb,mathtools}
\usepackage[scr=rsfs,cal=euler]{mathalfa}
\usepackage{xfrac}
\usepackage[unicode,hidelinks,bookmarks=false]{hyperref}
\newcommand{\ie}{i.e.\ }
\newcommand{\cf}{cf.\ }
\newcommand{\resp}{respectively\ }
\newcommand{\Subsection}{\subsection}
\providecommand{\nobreakdash}{\nobreak\hbox{-}}
\setlength{\emergencystretch}{2em}
\multlinegap0pt
\usepackage{bm}
\usepackage{bbm}
\usepackage{dsfont}
\usepackage{xspace}
\usepackage{cancel}
\usepackage{mathtools}
\usepackage{amsthm}
\makeatletter
\@ifundefined{theorem}{\newtheorem{theorem}{Theorem}}{}
\@ifundefined{assumptions}{\newtheorem{assumptions}[theorem]{Assumptions}}{}
\@ifundefined{proposition}{\newtheorem{proposition}[theorem]{Proposition}}{}
\makeatother
\title[Efficient high-resolution refinement in cryo-EM with stochas]{Efficient high-resolution refinement in cryo-EM with stochastic gradient descent}
\author{Bogdan Toader, Marcus A. Brubaker, Roy R. Lederman}
\date{Source: arXiv:2311.16100v3 (CC BY 4.0)}
\begin{document}
\maketitle

\noindent\emph{Source and licence.} Efficient high-resolution refinement in cryo-EM with stochastic gradient descent. Version arXiv:2311.16100v3 (CC BY 4.0), distributed under the Creative Commons Attribution 4.0 International licence, as recorded in the arXiv licence metadata for this exact version and shown on the abstract page https://arxiv.org/abs/2311.16100v3. Full rights notice: licenses/arxiv-2311.16100-CC-BY-4.0.txt.\par\smallskip

\noindent\textit{Editorial scope.} Counted unit: the Proposition whose source label is \texttt{prop:cond} in the article (source0.tex lines 660--675, source label \texttt{prop:cond}), delivered together with the article's own complete original proof of it (source0.tex lines 676--695, one \texttt{proof} environment of 20 lines). No mathematics is written, repaired or re-derived: statement and proof are the article's own text line for line. Required context retained uncounted: eq:cond H (lines 580--583); assump:projections (lines 626--635); eq:H diag (lines 650--653). Every definition, display and label of those lines is retained unchanged. Omitted from this package and not redistributed: 1--579, 584--625, 636--649, 654--659, 696--1520. Nothing inside a retained excerpt is omitted or altered: every retained body is delivered exactly as the article's lines are written, so no omission receipt of any kind accompanies this package. Citations: the retained excerpts carry no citation command at all, so no bibliography entry of the article is transcribed and no reference is redistributed. Reference adaptation: none was needed and none was made; every reference inside the delivered unit resolves inside it, so no unresolved reference or citation is produced. Printed slips kept literal: no printed word, symbol, hypothesis or number of the counted statement or of its proof was altered. Layout and class: the article's own class is \texttt{article} with packages including inputenc, fontenc, authblk, amsmath, amssymb, amsthm, commath, graphicx, caption, subcaption, float, algorithm, algpseudocode, makecell; the wrapper is the lane's pinned \texttt{amsart} editorial wrapper at 10pt on A4, which restates the theorem-like environments the retained text needs and adds the article's own macro definitions that the retained text needs (none), with extra wrapper packages (bm, bbm, dsfont, xspace, cancel, mathtools). The delivered numbering is the unit's own. Assets: no retained span contains a figure environment, a graphics-inclusion command, a PDF-page inclusion, a table or a TikZ picture; no image file is copied into this package, the package declares no asset dependency and no image file is redistributed with it. Licence: version arXiv:2311.16100v3 of this article is distributed under the Creative Commons Attribution 4.0 International licence, as recorded in the arXiv licence metadata for this exact version and shown on the abstract page https://arxiv.org/abs/2311.16100v3. A full rights notice is delivered at licenses/arxiv-2311.16100-CC-BY-4.0.txt. Characters: every retained excerpt is pure ASCII, so the delivered engine, XeLaTeX with its default Latin Modern text font, reports no missing character.
\begin{equation}
    \kappa(H) = \frac{\max_i H_{ii}}{\min_i H_{ii}}.
    \label{eq:cond H}    
\end{equation}

\begin{assumptions}\label{assump:projections}
\begin{enumerate}
    \item We assume that each voxel index $j$ is mapped at most once by the projection
        assignment function of an image $\Lambda_i$.
    \item Without loss of generality, we assume that $j=1$ is the index of the 
        voxel corresponding to the center of the coordinate axes. 
        Then, the voxel at $j=1$ is mapped by all projection operators $P_i$,
        $i=1,\ldots,N$, or equivalently, $\Omega_1 = N$.
\end{enumerate}
\end{assumptions}

\begin{equation}
    H_{jj} = |\Omega_j| + \lambda,
    \label{eq:H diag}
\end{equation}

\begin{proposition}[Condition number bound]\label{prop:cond}
    Let $M_x = M^2$ and $M_v = M^3 $ for grid length $M$ and let Assumptions~\ref{assump:projections} 
    hold for the nearest-neighbor interpolation projection matrices $P_i := P_{\phi_i}^{nn}$.
    Then, for any fixed number of images $N$, we have that:
    \begin{subequations}
    \begin{align}
        &\kappa(H) \geq \frac{N + \lambda}{N/M + \lambda},
        \qquad \forall M \leq N,
        \label{eq:cond number square lessthan}
        \\
        &\kappa(H) = \frac{N + \lambda}{\lambda},
        \quad\qquad \forall M > N.
    \end{align}
    \label{eq:cond number square}
    \end{subequations}
\end{proposition}

\begin{proof}
    Since the matrices $P_i$ are the nearest-neighbor interpolation matrices
    $P_{\phi_i}^{nn}$, the matrix $H$ is also diagonal and, according to~\eqref{eq:cond H},
    to compute $\kappa(H)$ we need to find the largest and smallest elements of $H$. 

    Using~\eqref{eq:H diag} and Assumptions~\ref{assump:projections}, we have that
    $\max_{j} H_{jj} = H_{11} + \lambda = N+\lambda$.

    To compute $\min_{j} H_{jj}$, note that the projection assignment functions
    $\Lambda_i$ for $i=1,\ldots,N$, map $NM^2$ image pixels to $M^3$ volume voxels.

    For $M > N$, there are more voxels than total pixels (in all the
    images), and so there exists a voxel $j^*$ such that $|\Omega_{j^*}| = 0$. 
    Then, $\min_{j}H_{jj} = H_{j^* j^*} = \lambda$, 
    and so $\kappa(H) = \frac{N+\lambda}{\lambda}$.
    For $M \leq N$, there are $NM^2$ pixels mapped to $M^3$ voxels, and therefore
    there exists a voxel $j^*$ such that $|\Omega_{j^*j^*}| \leq NM^2/M^3 = N/M$.
    Then, $\min_j H_{jj} \leq N/M + \lambda$, and so
    $\kappa(H) \geq \frac{N+\lambda}{N/M + \lambda}$.
\end{proof}

\end{document}
